\documentclass[10pt,a4paper]{amsart}
\usepackage[margin=19mm]{geometry}
\usepackage{amsmath,amssymb,mathtools}
\usepackage[scr=rsfs,cal=euler]{mathalfa}
\usepackage{xfrac}
\usepackage[unicode,hidelinks,bookmarks=false]{hyperref}
\newcommand{\ie}{i.e.\ }
\newcommand{\cf}{cf.\ }
\newcommand{\resp}{respectively\ }
\newcommand{\Subsection}{\subsection}
\providecommand{\nobreakdash}{\nobreak\hbox{-}}
\setlength{\emergencystretch}{2em}
\multlinegap0pt
\usepackage{amssymb}
\newtheorem{theorem}{Theorem}
\renewcommand{\P}{{\mathbb{P}}}
\title{Finiteness criteria for the solutions of a sequence of decomposable form inequalities}
\author{Si Duc Quang}
\date{Source: arXiv:2501.08607v2 (CC BY 4.0)}
\begin{document}
\maketitle

\noindent\emph{Source and licence.} Finiteness criteria for the solutions of a sequence of decomposable form inequalities. Version arXiv:2501.08607v2 (CC BY 4.0), distributed by arXiv under the Creative Commons Attribution 4.0 International licence, http://creativecommons.org/licenses/by/4.0/, as recorded in the arXiv licence metadata for this exact version and shown on the abstract page https://arxiv.org/abs/2501.08607v2. Full licence notice: \texttt{licenses/arxiv-2501.08607-CC-BY-4.0.txt}.\par\smallskip

\noindent\textit{Editorial scope.} Counted unit: the article's theorem 1.1 (source0.tex lines 105--114). The article's complete original proof of it is delivered with it (source0.tex lines 417--448, one \texttt{proof} environment of 32 lines, which the article places away from the statement). No mathematics is written, repaired or re-derived: the statement and the proof are the article's own lines, in the article's own notation, and no retained line is substituted or re-flowed.  Retained uncounted as the source-local context the proof needs: lines 222--229.  Nothing was omitted from any retained span, so the package carries no omission record.  No retained line contains a citation, so the wrapper carries no bibliography and no bibliography entry.  No retained line contains a figure environment, an image-inclusion command or drawing code, so no image is delivered and the package declares no asset dependency.  Editorial formatting only, in addition: an AMS \texttt{amsart} preamble that restates the article's own declarations and macros used by the retained lines, namely the environments \texttt{theorem} on separate counters, together with the standard packages those lines need.  The retained environments print in this document as Theorem 1 in order of appearance, whereas the article prints the same items with the numbering of its own full document.  The complete native source of the article as distributed in the arXiv e-print for this exact version is bundled unchanged as \texttt{sources/171615/source0.tex}.
\begin{theorem}\label{2.4} Let $k$ be a number field and $S \subset M_k$ be a finite set, containing all archimedean places. Further, let $\Lambda$ be an infinite set. Let $V$ be a subvariety of dimension $m$ of $\P^n(k)$ and $\mathcal Q=\{Q_1,\dots, Q_q\}$ be a family of moving hypersurfaces in $\P^n(k)$ indexed by $\Lambda$ with distributive constant $\Delta_{\mathcal Q,V}$ with respect to $V$. Let ${\bf x}=[x_0,\ldots,x_n]:\Lambda \to V$ be a sequence of points satisfying
$$h(Q_j(\alpha))=o(h({\bf x}(\alpha))) \text{ for all }\alpha\in\Lambda \text{ and }j=1, \dots,q$$
(i.e., for all $\delta>0$,  $h(Q_j(\alpha))\leq\delta h({\bf x}(\alpha))$ for all, but finitely many, $\alpha\in \Lambda).$
Then, for any $\epsilon >0,$ there exists an infinite index subset $A\subset \Lambda$ such that
\begin{eqnarray*} \sum_{v\in S}\sum_{j=1}^{q}\frac{\lambda_{Q_j(\alpha), v}({\bf x}(\alpha))}{\deg Q_j}\leqslant \left(\Delta_{\mathcal Q,V}\left(\frac{m}{2}+1\right)^2+\epsilon\right)h({\bf x}(\alpha)
)\end{eqnarray*}
holds for all $\alpha \in A.$
\end{theorem}

\begin{theorem}\label{1.1}
Let $\ell,m,q$ be positive integers and $\Delta\ge 1$ a positive constant. Let $k'$ be a finite extension of $k$ and $S \subset$ $M_{k}$ be a finite set containing all archimedean places. For $n=1,2, \ldots$, let $F_{n}({\bf x})=F_{n}(x_{0}, \ldots, x_{m}) \in \mathcal{O}_{S}[{\bf x}]$ be a sequence of semi-$q$-decomposable forms of degree $\ell$. Assume that $F_n=Q_{1,n} \cdots Q_{q,n}$ over $k'$ with $\deg Q_{j,n}=d_j, 1 \leqslant j \leqslant q$, for each $n$. Let $d=\max _{1 \leqslant j \leqslant q}d_j$. Assume that $\ell>d\Delta\left(\frac{m}{2}+1\right)^2$ and the family of $q$ moving hypersurfaces $\{Q_{1,n}, \ldots ,Q_{q,n}\}$ (indexed by $\mathbb N$) has distributive constant not exceeding $\Delta$ for each $n$. Let $c, \lambda$ be real numbers with $c>0, \lambda<\ell-d\Delta\left(\frac{m}{2}+1\right)^{2}$. Then, there does not exist an infinite sequence of $\mathcal{O}_{S}^{*}$-non-proportional $x_{n} \in \mathcal{O}_{S}^{m+1}, n=1,2, \ldots$, for which
\begin{align}\label{1.2}
0<\prod_{v \in S}\|F_{n}({\bf x}_{n})\|_{v} \leqslant c H_{S}^{\lambda}({\bf x}_{n})
\end{align}
and
\begin{align}\label{1.3}
h(F_{n})=o(h({\bf x}_{n})) \text { as } n\rightarrow \infty.
\end{align}
\end{theorem}

\begin{proof}[Proof of Theorem \ref{1.1}]
Let $S'$ be a subset of $M_{k'}$ which consists of the extension of the places of $S$ to $k'$. Then every $S$-integer in $k$ is also an $S'$-integer in $k'$. Moreover, we have $H_{S}({\bf x}_{n})=H_{S'}({\bf x}_{n})$ and
$$\prod_{v \in S}\|F_{n}({\bf x}_{n})\|_{v}=\prod_{w \in S'}\|F_{n}({\bf x}_{n})\|_{w} \text { for } {\bf x}_{n} \in \mathcal{O}_{S}^{m+1}.$$
So the inequality (\ref{1.2}) is preserved when we work on $k'$. Therefore, without loss of generality we may assume that $k'=k$. 

Suppose on the contrary that there is an infinite sequence ${\bf x}_{n}=(x_{0, n}, \ldots, x_{m,n}) \in \mathcal{O}_{S}^{m+1}$ which satisfies (\ref{1.2}) and (\ref{1.3}). If the heights $\{h({\bf x}_{n})\}_{n\geqslant 1}$ are bounded then the set $\{{\bf x}_{n};n\geqslant 1\}$ is actually a finite union of sets of proportional vectors. Therefore, we may suppose that $h({\bf x}_{n})$ are not bounded. 

Without loss of generality, we may assume that $h({\bf x}_{n})\rightarrow \infty$ as $n\rightarrow \infty$. Then, by the assumption, (\ref{1.3}) also holds. Furthermore it follows that $H_{S}({\bf x}_{n})\rightarrow \infty$ as $n\rightarrow \infty$ (since ${\bf x}_n\in\mathcal{O}_{S}^{m+1}$). By $\max _{0\leqslant j\leqslant m}h(Q_{j, n}) \leqslant h(F_{n})+O(1)$, from the inequality (\ref{1.3}) we deduce that 
$$\max _{0\leqslant j\leqslant m}h(Q_{j, n})=o(h({\bf x}_{n})) \text { as } n\rightarrow \infty.$$
Take $\epsilon$ a positive number small enough so that $\ell>d\Delta\left(\frac{m}{2}+1\right)^{2}+\lambda+\epsilon$.
By Theorem \ref{2.4}, there is an infinite subsequence ${\bf x}_{n_p}\in \mathcal{O}_{S}^{m+1},p=1,2, \ldots$, of $\{{\bf x}_{n}\}$, which without loss of generality we may assume to be $\{{\bf x}_{n}\}$ itself, such that
$$\sum_{v \in S} \sum_{j=1}^{q}\dfrac{1}{d_{j,n}} \log \frac{\|{\bf x}_{n}\|_{v}^{d_{j,n}} \cdot\|Q_{j, n}\|_{v}}{\|Q_{j, n}({\bf x}_{n})\|_{v}} \leqslant \left(\Delta\left(\frac{m}{2}+1\right)^{2}+\frac{\epsilon}{d}\right) h({\bf x}_{n}).$$
By the assumption that $d\geqslant \max_{1\leqslant j\leqslant n}d_{j,n}$, the above inequality implies that
\begin{align}\label{3.1}
\sum_{v \in S} \sum_{j=1}^{q}\log \frac{\|{\bf x}_{n}\|_{v}^{d_{j,n}}\cdot\|Q_{j, n}\|_{v}}{\|Q_{j, n}({\bf x}_{n})\|_{v}} \leqslant \left(d\Delta\left(\frac{m}{2}+1\right)^{2}+\epsilon\right) h({\bf x}_{n}).
\end{align}
On the other hand, since ${\bf x}_{n} \in \mathcal{O}_{S}^{m+1}$, we have
\begin{align}\label{3.2}
h({\bf x}_{n}) \leqslant \log H_{S}({\bf x}_{n}).
\end{align}
With the note that $\ell=\sum_{j=0}^md_{j,n}$ for each $n$, from (\ref{3.1}) and (\ref{3.2}) we have 
$$\prod_{v \in S} \frac{\|{\bf x}_{n}\|_{v}^{\ell} \cdot \prod_{j=1}^{q}\|Q_{j, n}\|_{v}}{\|F_{n}({\bf x}_{n})\|_{v}} \leq(H_{S}({\bf x}_{n}))^{d\Delta\left(\frac{m}{2}+1\right)^{2}+\epsilon}$$
whence
$$\frac{H_{S}^{\ell}({\bf x}_{n}) \cdot \prod_{v \in S} \prod_{j=1}^{q}\|Q_{j, n}\|_{v}}{\prod_{v \in S}\|F_{n}({\bf x}_{n})\|_{v}} \leq(H_{S}({\bf x}_{n}))^{d\Delta\left(\frac{m}{2}+1\right)^{2}+\epsilon}.$$
It is clear that $\prod_{v \in S} \prod_{j=1}^{q}\|Q_{j, n}\|_{v} \geq c' \prod_{v \in S}\|F_{n}\|_{v} \geq c'$ for some positive constant $c'$.
Therefore, from (\ref{1.2}), for $n=1,2, \ldots$, we have
\begin{align*}
H_{S}^{\ell}({\bf x}_{n})&\leqslant \frac{1}{c'} \prod_{v \in S}\|F_{n}({\bf x}_{n})\|_{v} \cdot(H_{S}({\bf x}_{n}))^{d\Delta\left(\frac{m}{2}+1\right)^{2}+\epsilon}\\ 
&\leqslant \frac{c}{c'}(H_{S}({\bf x}_{n}))^{d\Delta\left(\frac{m}{2}+1\right)^{2}+\epsilon+\lambda}. 
\end{align*}
Since $H_{S}({\bf x}_{n})\rightarrow \infty$ as $n\rightarrow \infty$, by letting $n\rightarrow \infty$ we get $\ell>d\Delta(\frac{m}{2}+1)^{2}+\lambda+\epsilon$. This is a contradiction. Hence, this completes the proof of Theorem \ref{1.1}.
\end{proof}

\end{document}
